\section{Дневник выполнения работы}
\begin{enumerate}
    \item Написание Python-скрипта для генерации файлов с $10^6$ командами (случайные, отсортированные, запросы поиска).
    \item Изучение утилиты \texttt{gcov}. Сборка программы с флагом \texttt{--coverage}. Анализ покрытия: логика вращений (Rotations) и балансировки успешно покрыта тестами.
    \item Изучение \texttt{Valgrind Memcheck} и \texttt{Massif}. Запуск программы, сбор данных о пиковом потреблении heap-памяти (68.83 MB) и отсутствие утечек.
    \item Компиляция с \texttt{-pg}, запуск и генерация отчётов \texttt{gprof}. Выявление узкого места: оверхед вызовов \texttt{std::unique\_ptr} и рекурсии.
    \item Переписывание архитектуры дерева: замена рекурсии на итеративные циклы с использованием \texttt{std::vector} в качестве стека пути, переход на сырые указатели.
    \item Повторное профилирование исправленной версии. Сбор метрик \texttt{time}, \texttt{gprof}, анализ улучшений. Подготовка отчёта.
\end{enumerate}

\section{Выводы}

Выполнение данной лабораторной работы продемонстрировало важность инструментального профилирования при разработке высоконагруженных систем. 

\textbf{Главный вывод:} алгоритмическая сложность $O(\log N)$ не является единственным фактором, определяющим скорость программы на практике. Инфраструктурный код языка (раскрутка рекурсивного стека вызовов, абстракции стандартной библиотеки вроде умных указателей) может занимать до 80\% процессорного времени.

Инструмент \texttt{gprof} позволил заглянуть внутрь процесса исполнения и увидеть то, чего не видно при статическом анализе кода. Инструмент \texttt{Valgrind Massif} позволил убедиться, что отказ от умных указателей в пользу сырых не ухудшил, а слегка улучшил профиль потребления оперативной памяти.

Итоговая (итеративная) реализация АВЛ-дерева оказалась быстрее, надёжнее (защищена от переполнения стека) и продемонстрировала способность конкурировать даже с хэш-таблицами (\texttt{std::unordered\_map}) на специфических (отсортированных) наборах данных. Я получил ценный опыт работы с профилировщиками в UNIX-системах, который буду применять в дальнейшей разработке на C/C++.
\pagebreak